\scanpage{74}
D'autre part, \deleted{\unsure{comme}} si $L$ est la partie linéaire de $G$, on a \deleted{$\unsure{H}$} $\underline{\unsure{\mathscr D|D}}_X(G)\simeq\underline{\mathscr D}_X(L)$, d'où
\[
H^i(\unsure{X|A}-a,G)\simeq H^i(X-a,L)\qquad(i\geq2,\text{ et }=\text{ on a vu pour }i\geq1)
\]
On sait par \unsure{là|l'a.} $H^i(X-a,G_m)=0$ puisque $X$ est factoriel, donc on aura $=$
\[
H^i(X-a,G)=H^i(X-a,U)\qquad(i\geq1)
\]
\deleted{\unsure{C'est}} D'ailleurs \unsure{on a}
\[
H^i(X-a,G_a)=0\qquad\text{si }i\neq\dim A-1
\]
donc
\[
H^i(X-a,U)=0\qquad\text{si }i\neq\dim A-1
\]
\unsure{Ceci achève de prouver ce qu'on veut|Ceci achève de prouver ce qui précède}.

\unsure{Noter aussi} $H_A^i(G)\xrightarrow{\sim}H_{\widehat A}^i(G)$ : rien ne change en \unsure{remplaçant} $A$ par $\widehat A$.
% archive sheet number: 73
